%%%%%%%%%%%%%%%%%%%%%%%%%%%%%%%%%%%%%%%%%%%%%%%%%%%%%%%%%%%%%%%%%%%%%%%%%%%%
% Executive Summary
%%%%%%%%%%%%%%%%%%%%%%%%%%%%%%%%%%%%%%%%%%%%%%%%%%%%%%%%%%%%%%%%%%%%%%%%%%%%

\chapter*{Executive Summary}
\addcontentsline{toc}{chapter}{Executive Summary}

Medora is our Spring final year project implementation of an AI-powered clinical intake and triage system. The project began in Fall 2025 as Healthify, a broad healthcare platform covering patient guidance, wearables, medical record summarization, and hospital resource optimization. In Spring 2026, we narrowed the work to the part we could implement deeply and evaluate honestly: structured intake, textbook-grounded triage, privacy-safe current evidence support, patient memory, doctor feedback, benchmarking, backend and frontend application infrastructure, full-stack agent integration, and hospital-server deployment preparation. A central goal is hospital-controlled deployment, where patient-facing inference runs on hospital infrastructure using validated open-source models so patient conversations and longitudinal memory do not have to leave the hospital environment.

The problem we address is the gap between unstructured patient complaints and the structured clinical information doctors need. Patients describe symptoms in everyday language, while clinicians need onset, duration, severity, associated symptoms, red flags, relevant history, and evidence-based differential reasoning. Medora is designed around that gap. It uses language models for conversation and structuring, but grounds clinical evidence in \textit{CURRENT Medical Diagnosis and Treatment 2022 (CMDT 2022)} and keeps every diagnosis provisional for physician review.

The implemented system has three technical layers. In the knowledge layer, we extracted CMDT 2022, corrected Private Use Area font corruption, kept chapter and section structure, generated 5,631 structure-aware chunks, embedded them with a medical sentence-transformer model, and stored them in ChromaDB. Retrieval validation on 20 clinical queries reached 90.0\% chapter Hit@1, 100.0\% Hit@3, and 100.0\% metadata filter pass rate. BGE cross-encoder reranking then improved the offline GPT-4o judge comparison score from 4.05 to 4.35 on a 1 to 5 scale.

In the agent and evidence layers, the Intake Agent uses 11 structured symptom objects to conduct a multi-turn interview, detect red flags, ask follow-ups for vague answers, support multi-symptom complaints, and produce a nine-section clinician handover note. The Triage Agent uses retrieval-augmented generation to produce a provisional diagnosis report with a primary diagnosis, differentials, reasoning, investigations, management considerations, red flags, and textbook sources. The Web Evidence Council adds a privacy-safe path for current public medical evidence: it removes patient identifiers, builds de-identified search queries, ranks trusted sources, extracts medical claims, reviews conflicts and safety, and returns physician-reviewable JSON. Patient Memory stores longitudinal clinical facts across sessions, while the Feedback Loop records doctor confirmation or rejection so reviewed cases can later become supervised training data.

The benchmark phase tested the RAG pipeline and web search fallback across textbook cases, MedQA-style cases, and rare MedCaseReasoning cases. On 50 textbook-derived cases, GPT-5.4-mini reached 74\% RAG accuracy, while the best open-source RAG run reported in the benchmark reached 42\%. On rare MedCaseReasoning cases, RAG accuracy dropped to 30\% because retrieval recall was only 22\%, while the web search path reached 42\% with GPT-5.4-mini by filling some textbook coverage gaps. Deployment preparation also began on an AWS EC2 g5.2xlarge instance with an NVIDIA A10G GPU, Ollama, and several open-source models.

The application layer turns the research pipeline into a working web prototype. The FastAPI backend uses PostgreSQL, SQLAlchemy, Alembic, JWT authentication, role guards, clinical report storage, RAG traceability tables, audit and performance log tables, and doctor feedback routes. The React and TypeScript frontend provides role-specific patient, doctor, and admin flows, including patient triage chat, protected routing, doctor report review, feedback submission, and admin user management. Phase 9 connects the real agent pipeline to the web application through an \texttt{AgentSessionManager}, shared model cache, background memory and feedback updates, and doctor-only diagnosis visibility.

The current implementation is a research prototype, not a finished medical product. It uses one primary textbook source, English-only interactions, prototype authentication and storage choices, and benchmark datasets rather than clinician-validated hospital cases. GPT-4o and GPT-5.4-family models were used for offline preprocessing, comparison, judging, benchmarking, and development-time testing, but they are not the intended hospital runtime models. A deployed hospital version must point the runtime LLM adapter to a validated hospital-hosted open-source or OpenAI-compatible local model, add production security hardening, validate the system clinically, and move active agent sessions to a production persistence strategy. Within these limits, Medora demonstrates a complete engineering path from raw medical text to a clinically structured, evidence-grounded, benchmarked, and web-accessible AI workflow.
